\section{Marco para evaluar productos de Agent}

Las secciones anteriores trataron por separado la experiencia, la variable del modelo, el contenedor, AgentRank y OS Agent. En esta sección esas ideas se organizan en un marco reutilizable: al evaluar una aplicación de Agent no basta con preguntar «¿qué tan potente es el modelo?»; hay que preguntar si define un nuevo entorno de tareas, si crea un nuevo grupo de usuarios, si hace que los token produzcan valor de forma más eficiente, si genera efectos de red y si cuenta con mecanismos de identidad y confianza.

\begin{knowledgebox}{Nota de clase: un producto de Agent no es un script de automatización}
Un script de automatización termina al completar una tarea; un producto de Agent debe generar de forma continua tareas, retroalimentación, invocaciones y relaciones dentro de un entorno. La ambición de YouWare no es generar un sitio web de una sola vez, sino lograr que las obras hechas con AI coding puedan publicarse, recibir interacción, ser invocadas e integrarse en una nueva red de creación.
\end{knowledgebox}

\noindent\begin{tabularx}{\textwidth}{p{0.20\textwidth}p{0.34\textwidth}X}
\toprule
Dimensión de evaluación & Pregunta que hay que hacer & Correspondencia en YouWare \\
\midrule
Nuevo grupo de usuarios & ¿Aparecen personas que antes no hacían esto? & Un Vibe Coder no equivale a un programador tradicional. \\
Contenedor & ¿Las obras tienen un nuevo entorno donde alojarse y compartirse? & El HTML/code se convierte en un sitio web accesible e interactivo. \\
Eficiencia de la inteligencia & ¿Los token se convierten de forma eficiente en obras y valor? & Pasar de conversar con el prompt a generar productos utilizables. \\
Efectos de red & ¿Usuarios, obras y Agent se refuerzan entre sí? & La distribución de obras, las invocaciones y AgentRank forman un círculo virtuoso. \\
Gobernanza de la confianza & ¿Hay identidad, permisos, auditoría y límites de seguridad? & Cuando los Agent se organicen en tribus, necesitarán una credencial de identidad y un candado con contraseña. \\
Dependencia del modelo & Cuando el modelo mejora, ¿el producto se beneficia de forma natural? & Confiar siempre en que el Model mejorará, pero que la aplicación construya el entorno. \\
\bottomrule
\end{tabularx}

\begin{importantbox}{Una buena aplicación de Agent debe avanzar en la misma dirección en que mejora el modelo}
Si cada avance del modelo obliga a reescribir una capa de andamiaje del producto, el producto está mal posicionado; si con cada avance del modelo la calidad de las obras, el alcance de las tareas y el valor para el usuario dentro del entorno crecen de forma natural, el producto está apostando por la variable correcta.
\end{importantbox}

\subsection{Tres errores frecuentes}

Esta subsección mira el marco desde el lado opuesto. El primer error es convertir el producto de Agent en una ventana de chat. Una ventana de chat es lo bastante general, pero el entorno es demasiado delgado: el usuario no sabe cómo convertir la inteligencia en obras que pueda compartir. El segundo error es convertir el producto de Agent en software empresarial tradicional. El software empresarial tiene ingresos claros, pero tiende a volverse pesado y reduce el espacio para nuevos grupos de usuarios y para la red del ecosistema. El tercer error es tratar la capacidad del modelo como barrera de entrada de la aplicación. Las empresas de modelos seguirán mejorando; la barrera de la aplicación tiene que estar en el entorno del usuario, el retorno de datos, los efectos de red y los mecanismos de confianza.

\begin{warningbox}{La barrera de una aplicación no puede depender solo de lanzar primero la UI}
La UI actual es fácil de copiar. Lo realmente difícil de copiar son las relaciones con los usuarios, la red de obras, las relaciones de invocación entre Agent, la identidad y la confianza, y la retroalimentación continua. Si nada de esto se consolida, llegar primero solo da visibilidad, no una barrera.
\end{warningbox}

\subsection{Resumen de la sección}

El marco para evaluar productos de Agent se resume en una frase: observar si crea un nuevo entorno de tareas y un nuevo grupo de usuarios, y si hace que la capacidad del modelo siga generando valor dentro de ese entorno. La apuesta central de YouWare es que el AI coding pasará de ser una capacidad de ingeniería a ser una capacidad creativa, y que la capacidad creativa necesita nuevos contenedores y nuevas redes.
